% !TEX root = main.tex
\subsection{Proving \SMS Function Evaluations}
\label{subsection:proving-sms-function-evaluations}

This subsection applies the PLONK framework of
Section~\ref{subsection:transforming-plonk-arithmetization-to-quadratic-constraints}
to the \SMS evaluation structure,
producing a single merged circuit $\CircuitMerge$
whose satisfiability the NIZKAoK in $\SMSSign$ proves.

\medskip
\noindent\textbf{Proving the hidden evaluation of $F$.}
Let $F$ be the signer's hidden policy
with $\nin^{(F)}$ inputs, $\ngate^{(F)}$ gates,
and $\nw^{(F)} \coloneqq \nin^{(F)} + 3\ngate^{(F)}$.
The signer evaluates $F$ on $\bx = (M \| \attr \| w)$
and retains the full computation trace.
Let $\plonkdescription_F$ denote $F$'s PLONK description
of the form~\eqref{equation:plonk-circuit-description}.
The signer then samples $\alpha, \beta \uniformly \FF$ from a random oracle
and assembles the witness vector:
\begin{equation}\label{equation:sms-ww-F}
\wit_F =
\bigl(
\plonkdescription_F \;\|\;
\bx \;\|\; \ba \;\|\; \bb \;\|\; \bc \;\|\;
(\hatw_i)_{i = 1}^{\ngate^{(F)}} \;\|\;
(\tildew_i)_{i = 1}^{\nw^{(F)}} \;\|\;
(w'_i)_{i = 1}^{\nw^{(F)}}
\bigr),
\end{equation}
where $\bx = (M \| \attr \| w)$
and all remaining components follow the definitions of
Section~\ref{subsection:transforming-plonk-arithmetization-to-quadratic-constraints}.
Viewing system~\eqref{equation:plonk-quadratic-constraints-with-challenges}
instantiated with $\plonkdescription_F$ as a single verification circuit
$\UniversalCircuitplonk_F$,
the signer must prove that $\wit_F$ satisfies $\UniversalCircuitplonk_F$
with public output $\out_F$.
Since $\UniversalCircuitplonk_F$ has
$\BigO(\nin^{(F)} + \ngate^{(F)})$ multiplications,
this proof is as efficient as proving the satisfiability
of a public circuit of the same size.

For the NIZKAoK to be sound, $\alpha$ and $\beta$ must be derived from a random oracle
after the prover commits to $\witpre$;
allowing the prover to choose $\alpha, \beta$ in advance would permit a cheating prover
to satisfy the grand product trivially.
This commitment-then-challenge structure is precisely the two-phase split
$(\witpre, \witsuf)$ of Section~\ref{subsection:transforming-plonk-arithmetization-to-quadratic-constraints},
and it drives the modification to ZKBoo in Section~\ref{section:overview-lattice-sms}.

\medskip
\noindent\textbf{Treating $G$ and $P$ as public circuits.}
Since $G$ and $P$ are public, PLONK's arithmetization is not needed for them.
They are treated directly as arithmetic circuits $\Circuit_G$ and $\Circuit_P$,
and their evaluations are proved as standard public-circuit satisfiability.

\medskip
\noindent\textbf{Merging into a single circuit.}
Given $\alpha$, $\beta$, $\wit_F$, $M$, $\attr$, $w$,
and intermediate outputs $\out_F$ and $\out_G$,
we merge $\UniversalCircuitplonk_F$, $\Circuit_G$, and $\Circuit_P$
into a single circuit $\CircuitMerge$ capturing:
$\alpha$, $\beta$, $\wit_F$, and $\out_F$ satisfy $\UniversalCircuitplonk_F$;
and $G(M, \attr, w) = \out_G$ and $P(\out_F, \out_G) = (\tM, \op)$.
A single NIZKAoK applied to $\CircuitMerge$ then certifies
the full evaluation $P(F(M, \attr, w), G(M, \attr, w))$
with public inputs $\alpha$, $\beta$, $(\tM, \op)$
and private witnesses $M$, $\attr$, $w$, $\wit_F$, $\out_F$, $\out_G$,
without leaking $F$'s circuit structure or computation trace.

After the Fiat--Shamir transform, $\alpha$ and $\beta$ are derived from
the random oracle and are publicly recomputable by the verifier;
the remaining witnesses $M$, $\attr$, $w$, $\wit_F$, $\out_F$, $\out_G$
are known only to the signer and hidden by the NIZKAoK.

\medskip
\noindent\textbf{Sizes.}
Let $\nadd^{(T)}$ and $\nmult^{(T)}$ denote the addition and multiplication gate counts
of function $T \in \{F, G, P\}$.
Circuit $\CircuitMerge$ has:
\begin{align*}
\text{additions:} \quad
&\BigO\bigl(\nin^{(F)} + \nadd^{(F)} + \nmult^{(F)}\bigr)
+ \nadd^{(G)} + \nadd^{(P)};
\\
\text{multiplications:} \quad
&\BigO\bigl(\nin^{(F)} + \nadd^{(F)} + \nmult^{(F)}\bigr)
+ \nmult^{(G)} + \nmult^{(P)}.
\end{align*}
For NIZKAoK systems whose proof size depends on the multiplication count,
the dominant cost is
$\BigO(\nin^{(F)} + \nadd^{(F)} + \nmult^{(F)}) + \nmult^{(G)} + \nmult^{(P)}$,
where $\nin^{(F)} = |M| + |\attr| + |w|$,
so the dominant multiplication count is
$\BigO(|M| + |\attr| + |w| + \nadd^{(F)} + \nmult^{(F)}) + \nmult^{(G)} + \nmult^{(P)}$,
which is quasilinear in the gate count of $F$ and linear in the multiplications of $G$ and $P$.
This confirms the proof-size claim of
Section~\ref{subsection:introduction-contributions}.
The ZKBoo-based instantiation exploits the two-phase split
$(\witpre, \witsuf)$ from
Section~\ref{subsection:transforming-plonk-arithmetization-to-quadratic-constraints}
to handle the random challenges $\alpha, \beta$;
details are in Section~\ref{section:overview-lattice-sms}.
